% ===================================================================
% MERGE PLAN: Files 03 (Architecture Patterns) and 04 (Modalities/Architectures)
% Date: 2026-01-27
% Purpose: Merge overlapping Results sections while preserving all content
% ===================================================================

% IDENTIFIED OVERLAPS:
%
% 1. ARCHITECTURE PATTERMS
%    - File 03: Mentions CNN, LSTM, ANN approaches with citations
%    - File 04: Repeats CNN, LSTM, ANN mentions with similar citations
%    - OVERLAP: Spahr (CNN ensemble), Elemam (CNNs), Wang (LSTM), Fine (ANN)
%
% 2. FORECASTING ARCHITECTURES
%    - File 03: Lists Meisel (10-unit LSTM), Nasseri (4-layer LSTM, 128 hidden)
%    - File 04: Repeats the same with slightly different wording
%    - OVERLAP: All LSTM forecasting studies
%
% 3. ENSEMBLE METHODS
%    - File 03: Spahr (30 CNNs), Dong (two-stage)
%    - File 04: Spahr mentioned again
%    - OVERLAP: Spahr ensemble details
%
% 4. HANDCRAFTED FEATURES
%    - File 03: Fine (594 features), Wang (attitude angles), Singh Rathore
%    - File 04: Fine (594 features, 100% sens), mentions 3/8 detection studies
%    - OVERLAP: Feature-based approach studies
%
% 5. ANOMALY DETECTION
%    - File 03: Reintjes, Ode
%    - File 04: Reintjes mentioned in validation section
%    - PARTIAL OVERLAP
%
% ===================================================================
% UNIQUE CONTENT IN EACH FILE:
% ===================================================================
%
% FILE 03 UNIQUE:
%    - Subsection "Learning Paradigms" with clear categorization
%    - End-to-end learning definition and examples
%    - Hybrid approaches (Vieluf, Elemam)
%    - Temporal Evolution (Borujeny 2013 as earliest, evolution trends)
%    - Window size ranges mentioned in key finding
%
% FILE 04 UNIQUE:
%    - Personalization Strategies (Global vs Patient-specific vs Mixed)
%    - Specific patient counts for each strategy (5 global, 4 patient-specific, 1 mixed)
%    - Detailed Global Models paragraph with Spahr, Fine, Dong, Elemam, Meisel
%    - Detailed Patient-Specific Models paragraph with Stirling, Nasseri, Ode, Singh Rathore
%    - Success rate comparisons (83-100% vs 62%)
%    - Cold start problem discussion
%    - Mixed Approaches (Vieluf)
%    - Validation Approaches (subject-split, temporal split)
%    - Computational Considerations (on-device vs cloud)
%    - Specific inference time (112 ms)
%
% ===================================================================
% PROPOSED MERGED STRUCTURE:
% ===================================================================
%
% \subsection{Architectures and Personalization}
% \label{sec:architecture-patterns}
%
%   \subsubsection{Learning Paradigms}
%     - From File 03 (well-structured categorization)
%     - Feature-based (4 studies): Fine, Wang, Singh Rathore
%     - End-to-end (2 studies, both forecasting): Meisel, Nasseri
%     - Hybrid (2 studies): Vieluf, Elemam
%     - Ensemble (2 studies): Spahr, Dong
%     - Anomaly detection (2 studies): Reintjes, Ode
%
%   \subsubsection{Architecture Patterns by Application}
%     - Merge content from both files
%     - CNN dominance in detection: Spahr, Elemam
%     - LSTM dominance in forecasting: Meisel, Nasseri, Stirling
%     - No forecasting uses CNN (temporal vs spatial patterns)
%     - Handcrafted features remain competitive (Fine 100% sens)
%     - Window size ranges
%
%   \subsubsection{Personalization Strategies}
%     - From File 04 (unique, important content)
%     - Global (5 studies): Spahr, Fine, Dong, Elemam, Meisel
%     - Patient-specific (4 studies): Stirling, Nasseri, Ode, Singh Rathore
%     - Mixed (1 study): Vieluf
%
%     \paragraph{Global Models.}
%       - On-device processing benefit
%       - Spahr 112ms inference
%       - Meisel LOSO: 62% patient success
%
%     \paragraph{Patient-Specific Models.}
%       - Higher success rates: 83-100%
%       - Data requirement: 14.6 months or 60+ days
%       - Cold start problem
%
%     \paragraph{Mixed Approaches.}
%       - Vieluf harmonic patterns + diary
%
%   \subsubsection{Validation Approaches}
%     - From File 04
%     - Subject-split: Reintjes (poor generalization)
%     - Temporal split: Nasseri
%     - LOSO: Meisel
%
%   \subsubsection{Temporal Evolution}
%     - From File 03 (unique content)
%     - Borujeny 2013 earliest
%     - 2020-2022: LSTM forecasting
%     - 2023-2026: Ensembles, hybrid, attention
%
%   \subsubsection{Computational Considerations}
%     - From File 04
%     - On-device: Spahr 112ms
%     - Cloud: Stirling, Nasseri
%     - Connectivity issues
%
%   \textbf{Key findings.}
%     - Detection: Architectural diversity (CNN, LSTM, ANN, anomaly)
%     - Forecasting: LSTM concentration, no CNNs
%     - Trade-off: Patient-specific (83-100% success, cold start) vs Global (62% LOSO, immediate deployment)
%
%   \textbf{Limitations.}
%     - Architecture description variability
%     - Detection studies rarely report patient-level success (2/8)
%
% ===================================================================
% CONTENT MAPPING:
% ===================================================================
%
% LEARNING PARADIGMS (keep from File 03, slight expansion):
%   - Fine: 594 features, 6-axis ACC+Gyro, ANN
%   - Wang: attitude angle features, LSTM 40 hidden
%   - Singh Rathore: multi-modal features, MLP
%   - Meisel: raw data, LSTM 10 units (E2E)
%   - Nasseri: FFT channels, LSTM 4-layer 128 hidden (E2E)
%   - Vieluf: DNN + harmonic features (hybrid)
%   - Elemam: CNN + rule-based fusion (hybrid - clarify classification)
%   - Spahr: 30 CNN ensemble, quantile aggregation
%   - Dong: two-stage, threshold pre-screening + CNN-LSTM
%   - Reintjes: anomaly detection, 3 methods, ECG
%   - Ode: self-attentive autoencoder, ECG RRI
%
% ARCHITECTURE PATTERNS (merged):
%   - CNN in detection: Spahr (96% sens, 0.0054/h), Elemam
%   - LSTM in forecasting: Meisel (10 units), Nasseri (128x4), Stirling (ensemble)
%   - No CNN in forecasting: temporal vs spatial
%   - Handcrafted features: Fine 100% sens, competitive
%   - Window sizes: 4s-5min (detection), 30s-60min (forecasting)
%
% PERSONALIZATION (keep from File 04):
%   - Global: Spahr (n=384), Fine (15 train, 3 test), Dong (n=68, 10-fold), Elemam (198 questionnaire, 30 HRV), Meisel (LOSO)
%   - Patient-specific: Stirling (n=11, weekly retraining), Nasseri (n=6), Ode (n=66, 99% confidence), Singh Rathore (n=1)
%   - Mixed: Vieluf (harmonic + diary)
%   - Success rates: 62% (global LOSO) vs 83-100% (patient-specific)
%   - Cold start: 14.6 months or 60+ days required
%
% VALIDATION (keep from File 04):
%   - Subject-split: Reintjes (FPR 0.11-65.62/h)
%   - Temporal split: Nasseri
%   - LOSO: Meisel
%
% COMPUTATIONAL (keep from File 04):
%   - On-device: Spahr 112ms, real-time
%   - Cloud: Stirling, Nasseri
%   - Smartphone: Stirling with weekly retraining
%
% TEMPORAL EVOLUTION (keep from File 03):
%   - 2013: Borujeny (ANN)
%   - 2020-2022: LSTM forecasting focus
%   - 2023-2026: Ensembles, hybrid, attention
%
% KEY FINDINGS (merged):
%   - Architecture diversity in detection vs LSTM concentration in forecasting
%   - Personalization trade-off: success vs cold start
%   - No CNN in forecasting
%
% LIMITATIONS (merged):
%   - Description variability (from 03)
%   - Patient-level reporting gap (from 04)
%
% ===================================================================
% ACTION ITEMS:
% ===================================================================
%
% 1. Remove duplicate mentions of studies between sections
% 2. Consolidate architecture descriptions (mention each study once in detail)
% 3. Keep Learning Paradigms structure from File 03
% 4. Keep Personalization section intact from File 04
% 5. Merge Architecture Patterns into single coherent subsection
% 6. Add Window Size discussion from File 04 to Architecture Patterns
% 7. Keep Temporal Evolution from File 03
% 8. Keep Validation and Computational from File 04
% 9. Consolidate Key Findings and Limitations
% 10. Update cross-references (remove "detailed in Section X" references)
% 11. Ensure all citation keys remain correct
% 12. Check line length and LaTeX formatting
%
% ===================================================================
